Efficient and reliable fault diagnosis of rotating machinery is critical for ensuring industrial safety and reducing maintenance costs. While deep learning has greatly advanced this field, existing methods often struggle with data scarcity and rely heavily on labor-intensive manual attribute labeling for zero-shot scenarios. To address these challenges, we propose the Cross-Modality Semantic Alignment Transformer (CMSA-Trans), a new framework that uses Large Language Model (LLM)-generated fault semantics for zero-shot fault diagnosis. CMSA-Trans employs a robust Time Series Transformer (TST) encoder to extract discriminative features from raw vibration signals and aligns them with semantic prototypes synthesized by an LLM. A cross-attention mechanism is integrated to enable deep semantic alignment between visual-like signal representations and textual descriptors, effectively capturing complex fault characteristics without requiring seen-class samples during training. Extensive experimental evaluations on the Case Western Reserve University (CWRU) and Southeast University (SEU) datasets show that CMSA-Trans greatly outperforms state-of-the-art zero-shot diagnosis methods. Our approach establishes a new approach for cross-modality diagnosis by bridging the gap between physical sensor data and high-level linguistic knowledge, offering a scalable solution for intelligent maintenance in dynamic industrial environments.

\begin{IEEEkeywords}
Zero-shot learning, fault diagnosis, Transformer, large language model, cross-modal alignment, rotating machinery, semantic prototyping.
\end{IEEEkeywords}

\begin{IEEEkeywords}
Zero-shot learning, fault diagnosis, Transformer, large language model, semantic alignment, rotating machinery.
\end{IEEEkeywords}
